\chapter*{Conclusions}

The work to develop MastoAnalyzer has been a fascinating and challenging journey into the world of the Fediverse and the challenges posed by detecting social bots. This research aimed to create a tool to monitor and understand interactions within the Fediverse, an area of the network still largely unexplored, tackling new complexities and discovering extraordinary opportunities.

Another crucial aspect was addressing the phenomenon of social bots, which can influence public discussions, spread misinformation, or even manipulate people's emotions. MastoAnalyzer demonstrates that by combining various analytical techniques, it is possible to detect bots with good accuracy, even in decentralized environments like the Fediverse.

Unfortunately, bots continue to evolve, becoming increasingly sophisticated and requiring continuous monitoring of ever-larger and more complex data volumes. However, the work done has laid a solid foundation to address these difficulties. MastoAnalyzer is designed to be a versatile and useful tool not only for researchers but also for the communities that populate the Fediverse. The results obtained show that this approach works and that we can improve the way we analyze decentralized platforms.

The structure of MastoAnalyzer is designed to efficiently and purposefully analyze data while respecting user privacy by collecting only publicly available data through APIs. This data is then processed using behavioral detection algorithms, all of which can be graphically visualized using the matplotlib library.

To better understand the dynamics between bots and human users in the context of the Fediverse, it was essential to analyze the data represented in the graphs produced by MastoAnalyzer. We can demonstrate that one of the most evident differences concerns behavior patterns: bots interact with greater regularity and repetitiveness, while human users are characterized by greater variability both temporally and in content. In fact, the graphs related to the temporal distribution of posts highlighted that bots often post at regular intervals, regardless of content or context, unlike human users who follow less predictable patterns based on time, events, or personal emotions.

Another aspect that characterizes bots is their tendency to form isolated clusters or predominantly interact with other bots, generating a less dense and more fragmented network. Human users, on the other hand, participate in more diverse conversations and contribute to more complex and interconnected networks. This data not only helps identify bots with greater precision but also highlights the importance of promoting authentic and dynamic networks within the Fediverse.

MastoAnalyzer is a modular tool, allowing new functionalities to be added or components to be replaced without affecting the entire system. It is also adaptable to other social networks within the Fediverse or even centralized ones.

In the future, we could enhance the capabilities of MastoAnalyzer by integrating more advanced technologies, such as more sophisticated artificial intelligence models, to further improve bot detection and expand the analysis to other aspects of online behavior, such as the impact of misinformation or the respect for privacy.

\addcontentsline{toc}{chapter}{Conclusions}


